\documentclass[12pt]{article}

% Font: Times New Roman
\usepackage{fontspec}
\setmainfont{Times New Roman}

% Margins: 1 inch all around
\usepackage[margin=1in]{geometry}

% Double spacing
\usepackage{setspace}
\doublespacing

% Paragraph formatting: indent first line, no extra space between paragraphs
\setlength{\parindent}{0.5in}
\setlength{\parskip}{0pt}

% Header
\usepackage{fancyhdr}
\pagestyle{fancy}
\fancyhf{}
\rhead{Aden \thepage}
\renewcommand{\headrulewidth}{0pt}

% For works cited
\usepackage{hanging}

% Disable hyphenation
\usepackage[none]{hyphenat}
\sloppy

% Prevent text from being cut off at page bottoms
\raggedbottom

\begin{document}

% Title block (MLA-ish style)
\thispagestyle{fancy}
{\setlength{\parindent}{0pt}
\noindent Jonah Aden\\
Professor Jayme R. Schlesinger\\
POLS UN1601: Introduction to International Politics\\
March 6, 2026

\begin{center}
\textbf{Assignment \#1}
\end{center}
}

\noindent \textbf{Question 1: Institutions and Anarchy}

The question of whether international institutions can substitute for a central authority and effectively replace anarchy is one of the central debates in international relations theory. Institutions clearly reduce the negative effects of anarchy by fostering cooperation, lowering uncertainty, and creating enforceable rules, but they do not and cannot serve as a substitute for central authority. This essay will discuss how liberal institutionalism demonstrates the ways institutions mitigate anarchy, and then examine the realist critique and empirical evidence that show institutions operate within anarchy, not as a replacement for it.

Liberal institutionalism provides the strongest case that institutions address the problems of anarchy. As Levy and Thompson argue, ``with proper domestic structures, international institutions, and economic interchange among themselves, states can mitigate the effects of international anarchy, minimize the violence-prone tendencies of the international system, promote further cooperation and reduce the frequency of war'' (Levy and Thompson 69). This shows that the liberal view sees institutions as a direct solution to the insecurity and conflict that anarchy produces. Institutions achieve this through specific mechanisms: they ``mitigate conflict, provide incentives to cooperate, reducing uncertainty in the international system'' (Schlesinger, ``Liberalism/EU'' lecture). Explicit institutions, such as the European Union, ``function by setting the rules of the game for participating members'' and ``encourage cooperation by lowering costs of cooperation and increasing the costs of conflict'' through express punishments for violations (Schlesinger, ``Liberalism/EU'' lecture). The EU represents the strongest empirical case, as supranational rules and enforcement mechanisms have sustained decades of peace among historically warring European states. Even Kant, one of the earliest advocates of institutional peace, envisioned a ``league of peace'' dedicated to ``the maintenance and security of the freedom of the state itself and of other states in league with it'' (Kant 8). This demonstrates that the idea of institutions as a way to manage anarchy has deep roots in political thought.

However, institutions fall critically short of replacing anarchy because they lack the enforcement power of true central authority. As the realist critique observes, ``power politics remain in institutions,'' which ``may provide an additional avenue through which major powers assert their influence'' over weaker states (Schlesinger, ``Liberalism/EU'' lecture). This shows that institutions can become tools of the powerful rather than neutral governance structures. When enforcement mechanisms are fragile, ``states may violate existing agreements'' (Schlesinger, ``Liberalism/EU'' lecture). If institutions truly replaced anarchy, states would not be able to simply leave, but Brexit showed that they can. The United Kingdom's withdrawal from the EU demonstrates the ``weakness of enforcement for explicit institutions'' and shows that membership remains voluntary even in the strongest institutional arrangements (Schlesinger, ``Liberalism/EU'' lecture). As Walt argues, ``There is no central authority that can protect states from one another,'' and cooperation ``is always somewhat fragile'' (Walt). Hobbes defined the anarchic condition as one in which ``men live without a common power to keep them all in awe'' (Hobbes, Ch. 13). Institutions, for all their benefits, do not constitute such a common power.

The crucial distinction is between governance \textit{within} anarchy and governance \textit{replacing} anarchy. Even Kant recognized this: his proposed federation ``would not have to be a state consisting of nations. That would be contradictory, since a state implies the relation of a superior (legislating) to an inferior (obeying)'' (Kant 6). Additionally, realist scholars argue that international institutions ``succeeded during the 1990s only because US dominance supported them,'' suggesting institutions depend on power distributions rather than transcending them (Eckstein). This shows that institutions are heavily reliant on the states that compose them, and the underlying logic of self-help and sovereignty persists.

In conclusion, institutions are a remarkable achievement in international politics. They make anarchy livable by creating rules, incentives, and norms that promote cooperation and reduce conflict. But they do not make anarchy obsolete. States retain sovereignty, can defect from agreements, and continue to operate under self-help logic. The question is not whether institutions help, because they clearly do, but whether they fundamentally alter the structure of international politics. They do not.

\newpage
\noindent \textbf{Question 2: Do Leaders Matter?}

Whether individual leaders shape international events or are shaped by them is a foundational question in international relations. Both perspectives carry strong evidence, and the most analytically useful answer is a conditional one: leaders matter most when structural constraints are ambiguous, mainly in crises, novel situations, and authoritarian contexts, but systemic pressures dominate when the external environment is compelling. This essay will analyze the evidence for both sides and argue that the influence of leaders depends on the interaction between individual agency and structural circumstance.

Individual-level theories provide compelling evidence that leaders are decisive in certain conditions. As Levy and Thompson argue, variations in leaders' beliefs ``arise from differences in political socialization, personality, education, formative experiences,'' and ``when individual beliefs vary, different decision-makers will respond differently under similar situations.'' As Hermann et al. emphasize, ``who leads matters'' (Levy and Thompson 104--105). This shows that the individual at the helm of a state can drastically alter the course of its foreign policy. The case of Nixon's opening to China powerfully illustrates this: ``having established a strong and persistent anti-Communist record during the 1950's and 1960's, President Nixon then opened the door to the international legitimization of the People's Republic of China'' (Cukierman and Tommasi). Nixon's individual identity and credibility enabled a policy shift that no other leader could have pursued, and structure alone cannot explain this outcome. Additionally, Carter argues that Xi Jinping's personal political survival drives China's foreign policy, contending that ``China's domestic politics compel risky behavior on the international stage'' through diversionary aggression (Carter). In the case of the U.S. withdrawal from Afghanistan, both the Trump and Biden administrations were ``held captive'' by campaign slogans, ``allowing political considerations to override intelligence assessments,'' and the decision ``was predetermined regardless of warnings about alarming outcomes'' (London). This shows that in each of these cases the leader's individual characteristics and personal political concerns drove outcomes that structural analysis alone cannot predict.

Yet when the external environment is compelling, individual differences wash out. Arnold Wolfers' analogy captures this: ``if the house is on fire we do not have to know anything about the individuals in it to predict that they will rush for the exits'' (Jervis). So, even though leaders clearly matter in ambiguous situations, the opposite is also true when the environment is constraining. Jervis finds that across post-Cold War American presidents, ``it would be hard to infer when a new person entered the Oval Office,'' and despite vastly different personalities, policy remained strikingly continuous (Jervis). The organizational process model explains why: organizations ``select which of their pre-existing routines best fits the situation'' and implement those routines, meaning ``the best predictor of how an organization behaves at time t is how they behaved at time t-1'' (Levy and Thompson 137). Moreover, ``political leaders do not implement their own policies, and there is often `slippage' between the foreign policy decisions of political leaders and the ways in which their decisions are actually implemented'' (Levy and Thompson 146). This shows that bureaucratic inertia, organizational routines, and systemic pressures create continuity that persists regardless of who occupies office.

The resolution lies in recognizing the interaction between individual and structural factors. As Jervis writes, ``not all situations are created equal. Not all people are either, and where one person might feel trapped another could see room for choice. There is an interaction between individuating and circumstantial factors'' (Jervis). This shows that the external environment sets the range of possible outcomes and the individual leader determines the specific outcome within that range, meaning the question of whether leaders matter depends greatly on the situation they find themselves in.

In conclusion, when the situation is ambiguous leaders can be decisive, as we can see with Nixon's credibility enabling the opening to China, Xi's political survival driving diversionary aggression, and two administrations allowing campaign promises to override intelligence on Afghanistan. However, when the external environment is compelling, structure dominates and anyone in office would likely act the same way. Understanding this interaction between leaders and structure is more analytically useful than a blanket yes or no answer to the question of whether leaders matter.

\newpage
\noindent \textbf{Question 3: Threat Perception and Ukraine}

The fact that North Korean nuclear weapons constitute a different threat to the United States than British nuclear weapons tells us something fundamentally important about how international politics works. Threats are not determined by material capability alone, they are socially constructed through identity and shared ideas, and this is the central argument of constructivism. This essay will first explain the constructivist principle behind this observation, and then apply the same logic to argue that Russia's invasion of Ukraine is best understood as a response to the ideational threat posed by Ukraine's shift toward Western democratic identity and NATO membership.

The nuclear example illustrates the constructivist argument that material power alone does not determine threat. As Wendt's famous example demonstrates, ``500 British nuclear weapons are less threatening to the United States than five North Korean nuclear weapons.'' The difference ``lies not in material weapons but in the meanings assigned to them, the ideational structures that matter more than material ones'' (Theys). The course lecture reinforces this directly: ``Nuclear capabilities in the United Kingdom have different meaning to the United States than nuclear capabilities in North Korea'' because ``ideas shape perception of states' intentions'' (Schlesinger, ``Constructivism'' lecture). Constructivism holds that politics is ``historically and socially constructed'' and ``emphasizes ideational factors over material factors'' (Schlesinger, ``Constructivism'' lecture). The United Kingdom is a NATO ally with a shared democratic identity and common values with the United States, while North Korea is an adversary with a hostile identity and opposing ideology. Even though the weapons are materially similar, the identities of the states holding them produce entirely different threat perceptions. This shows that identity constitutes interest and interest determines threat. As Wendt argues, ``anarchy is what states make of it,'' meaning states interpret the international system through the meanings they assign, and ``states possess multiple socially constructed identities that signal their interests'' (Theys).

This constructivist logic applies directly to Russia's invasion of Ukraine. Ukraine's application for NATO membership carried profound identity implications because it signaled a shift in national identity from a post-Soviet state within Russia's sphere to a Western democracy aligned with NATO. As the course lecture on NATO and Russia-Ukraine notes, ``NATO and its member states signal intentions through shared identity'' and ``membership (non-membership) status may shape perceptions of states' interests'' (Schlesinger, ``NATO/Russia-Ukraine'' lecture). So, the constructivist explanation emphasizes ``identity (allies/adversaries) implications of Ukraine's application for NATO membership, national identity (Russian/Ukrainian), status as a post-Soviet state'' and ``signaling domestic politics under democratic governance'' (Schlesinger, ``NATO/Russia-Ukraine'' lecture). NATO functions as an ``ideological balancer,'' and its enlargement activates ``Russian-Ukrainian historically constructed identities, populations, cultures, languages, territories'' (Schlesinger, ``NATO/Russia-Ukraine'' lecture). Walt captures the constructivist perspective on NATO expansion: it incorporates states ``within the Western security community, whose members share a common identity that has made war largely unthinkable'' (Walt, ``One World, Many Theories''). For Russia, Ukraine joining this community would redefine Ukraine's identity and, by extension, the meaning of the entire post-Soviet space. This shows that in order to understand Russia's response, we have to look beyond the material balance of power and examine the identities at play.

While realism identifies NATO enlargement as altering ``the power dynamics between Russia and NATO-member states'' (Schlesinger, ``NATO/Russia-Ukraine'' lecture), it cannot fully explain why this shift is perceived as threatening in the first place. Does the same military alliance carry different meaning depending on who joins it? Constructivism says yes. As constructivists demonstrate, ``state behavior reflects not only material power distribution and geography but also ideas, identities, and norms,'' and ``international reality isn't fixed, it remains subject to change through actors' reconceptualization of meaning and identity'' (Theys). Russia invaded Ukraine because Ukraine's identity transformation threatened the very meaning Russia assigns to its role as a great power and its place in the post-Soviet space, which goes beyond a simple material power shift.

In conclusion, the same logic that explains why British and North Korean nuclear weapons carry different meaning also explains why NATO expansion is so threatening to Russia. The same weapons and the same alliance produce different threat perceptions depending on the identities involved, and Russia's invasion of Ukraine was driven by what NATO membership signifies for Ukraine's identity and its place in the post-Soviet space. Constructivism shows us that the international system is built on ideas and identities as much as it is built on material power.

\newpage
\begin{center}
\textbf{Works Cited}
\end{center}

\begin{hangparas}{0.5in}{1}

Carter, Erin Baggott. ``The Diversionary Aggression of Authoritarian Leaders.'' Lecture reading, POLS UN1601, Columbia University, 2019.

Cukierman, Alex, and Mariano Tommasi. ``When Does It Take a Nixon to Go to China?'' \textit{The American Economic Review}, vol. 88, no. 1, 1998, pp. 180--197.

Eckstein, Arthur. ``Anarchy in International Relations Theory: The Neorealist-Neoliberal Debate.'' Lecture reading, POLS UN1601, Columbia University, 2021.

Hobbes, Thomas. \textit{Leviathan}. 1651.

Jervis, Robert. ``Do Leaders Matter and How Would We Know?'' \textit{Security Studies}, vol. 22, no. 2, 2013, pp. 153--179.

Kant, Immanuel. ``Perpetual Peace: A Philosophical Sketch.'' 1795.

Levy, Jack S., and William R. Thompson. \textit{Causes of War}. Wiley-Blackwell, 2010.

London, Douglas. ``CIA's Former Counterterrorism Chief for the Region: Afghanistan, Not an Intelligence Failure --- Something Much Worse.'' Lecture reading, POLS UN1601, Columbia University, 2021.

Schlesinger, Jayme R. ``Constructivism.'' Lecture, POLS UN1601: Introduction to International Politics, Columbia University, Spring 2026.

Schlesinger, Jayme R. ``Liberalism / Case Study: European Union.'' Lecture, POLS UN1601: Introduction to International Politics, Columbia University, Spring 2026.

Schlesinger, Jayme R. ``NATO / Russia-Ukraine War.'' Lecture, POLS UN1601: Introduction to International Politics, Columbia University, Spring 2026.

Theys, Sarina. ``Introducing Constructivism in International Relations Theory.'' \textit{E-International Relations}, 2018.

Walt, Stephen M. ``One World, Many Theories.'' \textit{Foreign Policy}, no. 110, 1998, pp. 29--46.

Walt, Stephen M. ``The World Wants You to Think Like a Realist.'' \textit{Foreign Policy}, 2018.

\end{hangparas}

\end{document}
